\section{Vocales largas}
\textbf{Qué son:} Vocales con duración y apertura mayores, se escriben como vocal doble o vocal + consonante simple: \textit{aa, ee, oo, uu}.

\textbf{Por qué importan:} Cambian significado frente a la versión corta y afectan el número de sílabas.

\textbf{Cómo practicarlas:}
\begin{itemize}[leftmargin=*]
  \item Pares mínimos: \textit{maan/man}, \textit{been/ben}, \textit{boon/bon}, \textit{huur/hur?} (contraste \textit{uu/u}).
  \item Lee en voz alta listas: \textit{kaas, maan, meer, boom, muur}.
  \item Marca duración: mantén la vocal dos tiempos, consonante corta.
\end{itemize}
